\subsection{Physikalische Re/Ra-Zuordnung}\label{sec:eabc-re-ra}

Vereinigter EABC-Zeugen-Katalog am Kepler-Anker $M=113160$:
96 rationale Fixpunkte (Radius~4) und
9 kanonische $\Phi$-Zeugen (Radien~2/3, dedupliziert nach Schale $\alpha+\beta\Phi$).

\paragraph{Kalibrierung Reynolds ($a=0$).}
Im rationalen Sektor gilt $M/\mathrm{Norm}\in\mathbb{Z}$.
Die reziproken Normen $1/\mathrm{Norm}$ koppeln mit der $M$-Kaskade
über
\[
  \kappa_{Re} = \frac{M}{Re_{\mathrm{krit}}}, \qquad
  Re = \frac{M/\mathrm{Norm}}{\kappa_{Re}}
     = \frac{Re_{\mathrm{krit}}}{\mathrm{Norm}}.
\]
Am Einheitsschalenpunkt ($\mathrm{Norm}=1$, $M/\mathrm{Norm}=M$) emergiert
$Re_{\mathrm{krit}} = 2300$ exakt
(berechnet: 2300.000000, Abweichung $0.00e+00\,\%$).
Kaskade: $113160/56580=2$, $56580/37720=3/2$; strukturell $M=3\cdot 37720$.

\begin{table}[ht]
  \centering
  \caption{Rationale Norm-Schalen und $M$-Kaskade (Radius~4).}
  \label{tab:eabc-re-schalen}
  \begin{tabular}{rrrrr}
    \toprule
    $\mathrm{Norm}$ & $M/\mathrm{Norm}$ & $1/\mathrm{Norm}$ & Anzahl & Kaskade \\
    \midrule
    1 & 113160 & 1.0000 & 6 & --- \\
    2 & 56580 & 0.5000 & 6 & 2 \\
    3 & 37720 & 0.3333 & 30 & 1.5 \\
    4 & 28290 & 0.2500 & 6 & 1.333 \\
    5 & 22632 & 0.2000 & 24 & 1.25 \\
    6 & 18860 & 0.1667 & 24 & 1.2 \\
    \bottomrule
  \end{tabular}
\end{table}

\begin{table}[ht]
  \centering
  \caption{Auszug Re-Zuordnung (rationaler Sektor).}
  \label{tab:eabc-re-zuordnung}
  \begin{tabular}{rrrrr}
    \toprule
    $(c_e,c_a,c_b,c_c)$ & $\mathrm{Norm}$ & $M/\mathrm{Norm}$ & $1/\mathrm{Norm}$ & $Re$ \\
    \midrule
    (-1, 0, 0, 0) & 1.0 & 113160 & 1.0000 & 2300.0 \\
    (0, 0, -1, 0) & 1.0 & 113160 & 1.0000 & 2300.0 \\
    (0, 0, 0, -1) & 1.0 & 113160 & 1.0000 & 2300.0 \\
    (0, 0, 0, 1) & 1.0 & 113160 & 1.0000 & 2300.0 \\
    (0, 0, 1, 0) & 1.0 & 113160 & 1.0000 & 2300.0 \\
    (1, 0, 0, 0) & 1.0 & 113160 & 1.0000 & 2300.0 \\
    (-2, 0, 0, 0) & 2.0 & 56580 & 0.5000 & 1150.0 \\
    (0, 0, -2, 0) & 2.0 & 56580 & 0.5000 & 1150.0 \\
    (0, 0, 0, -2) & 2.0 & 56580 & 0.5000 & 1150.0 \\
    (0, 0, 0, 2) & 2.0 & 56580 & 0.5000 & 1150.0 \\
    (0, 0, 2, 0) & 2.0 & 56580 & 0.5000 & 1150.0 \\
    (2, 0, 0, 0) & 2.0 & 56580 & 0.5000 & 1150.0 \\
    \bottomrule
  \end{tabular}
\end{table}

\paragraph{Kalibrierung Rayleigh ($a\neq 0$).}
Im $\Phi$-Sektor gilt $M^2/\mathrm{Norm}^2\in\mathbb{Z}[\Phi]$.
Der Urzeuge $(0,-1,0,0)$ mit $\mathrm{Norm}=\Phi$ und Feldnorm
$N(1+\Phi)=1$ kalibriert
\[
  \kappa_{Ra} = \frac{Ra_{\mathrm{krit}}}{N(\alpha+\beta\Phi)},
  \qquad Ra = N(\alpha+\beta\Phi)\cdot Ra_{\mathrm{krit}}.
\]
Am ikosaedrischen Anker emergiert $Ra_{\mathrm{krit}} = 1708$ exakt
(berechnet: 1708.000000, Abweichung $0.00e+00\,\%$).

\begin{table}[ht]
  \centering
  \caption{$\Phi$-Sektor: Rayleigh-Zuordnung (kanonische Zeugen).}
  \label{tab:eabc-ra-zuordnung}
  \begin{tabular}{rrrrrr}
    \toprule
    $(c_e,c_a,c_b,c_c)$ & $\mathrm{Norm}$ & $N$ & $M^2/\mathrm{Norm}^2$ & $Ra$ & Urz. \\
    \midrule
    (0, -1, 0, 0) & 1.6180 & 1 & $25610371200 - 12805185600\Phi$ & 1708.0 & \checkmark \\
    (-1, -1, 0, 0) & 1.9021 & 5 & $7683111360 - 2561037120\Phi$ & 8540.0 &  \\
    (-2, -1, -1, 0) & 2.7601 & 41 & $2186251200 - 312321600\Phi$ & 70028.0 &  \\
    (0, -2, 0, 0) & 3.2361 & 16 & $6402592800 - 3201296400\Phi$ & 27328.0 &  \\
    (-2, -2, 0, 0) & 3.8042 & 80 & $1920777840 - 640259280\Phi$ & 136640.0 &  \\
    (-2, -2, -2, -1) & 4.4127 & 205 & $1061893440 - 249857280\Phi$ & 350140.0 &  \\
    (0, -3, 0, 0) & 4.8541 & 81 & $2845596800 - 1422798400\Phi$ & 138348.0 &  \\
    (-2, -3, 0, 0) & 5.2500 & 205 & $1374215040 - 562178880\Phi$ & 350140.0 &  \\
    (-3, -3, 0, 0) & 5.7063 & 405 & $853679040 - 284559680\Phi$ & 691740.0 &  \\
    \bottomrule
  \end{tabular}
\end{table}

\noindent\textbf{Skalenfaktoren:}
$\kappa_{Re} = 49.200000$,
$\kappa_{Ra} = 1708.000000$.
Deterministischer Lauf: \texttt{eabc\_re\_ra\_zuordnung.py}.
